\section{Representation dominance for reflection-only sampling}
\label{sec:representation-sampler-dominance}

The following theorem separates an implementation cost statement from an
estimator comparison. Let $A_i$ be a completely clean financial
preparation with its adjoint, and let $S_{0,i}$ reflect the complete
invariant input-zero subspace. The good reflection is a single Clifford
$Z$ on the already computed flag. Write $a_i=T(A_i)$ and
$s_i=T(S_{0,i})$, counting $T$ and $T^\dagger$ together. The complete
adjoint costs $a_i$. Measurement, classical conjunction of the observed
input-zero bits and conditional Pauli resets require no $T$ gates under
the declared logical architecture; their non-$T$ counts and allocated
wire counts remain separately charged. Scratch may be excluded from
the zero reflection only after its cleanliness invariant has been
proved. This theorem assumes those preparations and reflections, and
does not supply a missing coherent state-access or cleanliness proof.
The declared charge adds the complete retained preparation, adjoint and
reflection templates. It applies no cancellation or optimisation across
their boundaries. A different composition or cross-template optimisation
must be recosted before applying the theorem; the identities are not a
lower bound on every circuit realisation of the same polynomial signal.

\begin{theorem}[All-degree reflection-sampler dominance]
\label{thm:representation-sampler-dominance}
For every integer $d\geq1$, the retained parity-specific reflection-only
sampler has complete per-shot logical $T$ count
\begin{equation}
 t_i(d)=
 \begin{cases}
   da_i+(d-1)s_i/2,&d\text{ odd},\\
   da_i+(d/2-1)s_i,&d\text{ even}.
 \end{cases}
 \label{eq:representation-parity-cost}
\end{equation}
Candidate $c$ has $t_c(d)\leq t_b(d)$ for every positive integer $d$
if and only if
\begin{equation}
 a_c\leq a_b,\qquad 2a_c+s_c\leq2a_b+s_b.
 \label{eq:representation-all-degree-condition}
\end{equation}
If $a_b>0$ and $2a_b+s_b>0$, then every degree obeys the uniform ratio
bound
\begin{equation}
 \frac{t_c(d)}{t_b(d)}\leq
 \max\left\{\frac{a_c}{a_b},
             \frac{2a_c+s_c}{2a_b+s_b}\right\}.
 \label{eq:representation-all-degree-ratio}
\end{equation}
Both endpoint ratios are necessary for a uniform bound of this form:
the first is attained at $d=1,2$ and the second is approached as
$d\longrightarrow\infty$. Strict superiority for every degree holds
if and only if $a_c<a_b$ and $2a_c+s_c\leq2a_b+s_b$.
\end{theorem}
\begin{proof}
Put $Q_i=-A_iS_{0,i}A_i^\dagger S_g$. The unconditional global minus
can be omitted because it changes a terminal shot only by a global
phase. It cannot be omitted from a controlled-QPE comparison. For
$d=2k+1$, the chronological sampler prepares $A_i$ and then applies
$Q_i$ exactly $k$ times. It uses $k+1$ preparations, $k$ adjoints,
$k$ zero reflections and $k$ Clifford good reflections. For $d=2k$
with $k\geq1$, it prepares $A_i$, applies $Q_i^{k-1}$, applies $S_g$
and then $A_i^\dagger$, observing the invariant input-zero event. It
uses $k$ preparations, $k$ adjoints, $k-1$ zero reflections and $k$
Clifford good reflections. Complete measurement and reset are appended
to both constructions. Their declared $T$ charge is zero, so these
counts prove \eqref{eq:representation-parity-cost}.

Let $q_i=2a_i+s_i$, $\delta a=a_c-a_b$ and $\delta q=q_c-q_b$.
The odd costs are $a_i+kq_i$ for $k\geq0$ and the even costs are
$2a_i+(k-1)q_i$ for $k\geq1$. If $\delta a\leq0$ and
$\delta q\leq0$, subtracting these formulae gives nonpositive
differences in both parities. Conversely, degree one requires
$\delta a\leq0$. If $\delta q>0$, the odd difference
$\delta a+k\delta q$ is positive at a sufficiently large integer
$k$, contradicting all-degree dominance. This proves necessity and
sufficiency. At $d=1$ strict dominance requires $\delta a<0$; together
with $\delta q\leq0$ this is also sufficient in both parities.

For odd degrees the ratio of costs is
$(a_c+kq_c)/(a_b+kq_b)$. It is a convex combination of $a_c/a_b$
and $q_c/q_b$, with nonnegative weights
$a_b/(a_b+kq_b)$ and $kq_b/(a_b+kq_b)$. For even degrees replace
$a_i$ by $2a_i$ and $k$ by $k-1$; the two endpoint ratios are
unchanged. A convex combination is bounded by the larger endpoint.
Degrees one and two attain $a_c/a_b$. Along either parity the ratio
tends to $q_c/q_b$ as its degree tends to infinity, proving the final
uniform sharpness statement.
\end{proof}

\begin{corollary}[Same-history adaptive sampler comparison]
\label{cor:representation-same-history}
Suppose two clean representations have the same implemented good
probability $p$, and an estimator uses the same initial prior, tolerance,
confidence allocation and deterministic history-dependent degree/batch
policy for both. Under ideal logical operation their complete reflection
shots can be coupled to have identical observed Bernoulli outcomes,
chosen degrees, batch sizes and stopping decisions. If
\eqref{eq:representation-all-degree-condition} holds, every coupled
history has no larger candidate reflection-shot $T$ cost. If
\eqref{eq:representation-all-degree-ratio} has right side $r$, then
the total candidate reflection-shot cost on every history is at most
$r$ times the corresponding rival reflection-shot cost. Consequently
the same comparison holds for expectations whenever they exist.
\end{corollary}
\begin{proof}
The two-dimensional reflection algebra gives the same
Chebyshev Bernoulli parameter $q_d(p)$ for both clean representations
at each degree. This includes $p=0,1$ by direct endpoint evaluation,
without normalising a zero-probability basis component. At each shot,
use a fresh common uniform random variable $V$ to define both observed
bits as $\mathbf1_{\{V\leq q_d(p)\}}$. Inductively, identical preceding
histories imply the same next degree and batch size, which imply
identical next observation under this coupling. The deterministic
policy then makes the same next decision. Apply the pointwise
per-shot inequality to each degree in the common history and sum its
nonnegative costs. Taking expectation preserves the inequality.
For unbounded stopping histories, apply the same argument to finite
prefixes and then use monotone convergence. Any additional fallback
or non-reflection primitive requires its own cost comparison; it is not
included by this corollary merely because it belongs to the estimator.
\end{proof}

The continuous ideal numerical, raw-payoff and shared-threshold
constructions target the same financial mean and hence the same ideal
Grover angle after a matched monetary normaliser. Their separately
rounded finite implementations need not have identical $p$. Their
deterministic allowances, probability priors, admissible degrees and
stopping histories can therefore differ. The all-degree cost theorem
compares the actual complete preparation resources at matched degrees;
the same-history corollary is conditional on matched implemented
probability and controller information. Neither theorem alone proves
smaller realised adaptive costs, a new amplitude-estimation query law,
dominance over non-reflection estimators or physical quantum advantage.
The generic reflection algebra is established. The project-specific
positive consequence is an exact, fully charged representation
comparison if the retained preparations satisfy its two affine
inequalities; mathematical priority remains a separate literature
question.
